برای نشان دادن اینکه یک ماشین تورینگ قطعی دوبعدی با یک ماشین تورینگ تک‌نواره استاندارد معادل است، باید ثابت کنیم که هر کدام می‌توانند دیگری را شبیه‌سازی کنند.

\subsubsection*{شبیه‌سازی ماشین تورینگ استاندارد توسط ماشین تورینگ دوبعدی}
این حالت بسیار ساده است. ماشین تورینگ دوبعدی می‌تواند به سادگی نوار افقی اولیه را به عنوان تنها نوار خود در نظر بگیرد و هرگز حرکات بالا و پایین را انجام ندهد. با محدود کردن حرکات به چپ، راست و ثابت ماندن، این ماشین دقیقا مانند یک ماشین تورینگ تک‌نواره استاندارد عمل خواهد کرد.

\subsubsection*{شبیه‌سازی ماشین تورینگ دوبعدی توسط ماشین تورینگ تک‌نواره}
برای این کار، ماشین تک‌نواره باید صفحه دوبعدی نامتناهی را روی یک نوار یک‌بعدی مدل‌سازی کند. ایده اصلی این است که فقط خانه‌هایی از صفحه دوبعدی که محتوای آن‌ها غیرخالی است را همراه با مختصاتشان روی نوار ذخیره کنیم.

\begin{itemize}
    \item {نحوه ذخیره‌سازی روی نوار:} 

    یک دستگاه مختصات دکارتی در نظر می‌گیریم که موقعیت اولیه هد در مختصات $(0,0)$ قرار دارد. روی نوار ماشین یک‌بعدی، محتوای هر خانه غیرخالی از صفحه دوبعدی را به صورت یک رشته به فرم $(x, y, a)$ ذخیره می‌کنیم که در آن $x$ و $y$ مختصات خانه و $a$ نماد موجود در آن است. این بلاک‌ها با یک نماد جداکننده مانند $\#$ از هم جدا می‌شوند.
    
    \item {ذخیره وضعیت فعلی:}

    ماشین یک‌بعدی مختصات فعلی هد ماشین دوبعدی یعنی $(x_c, y_c)$ را در یک شیار جداگانه روی نوار یا در ابتدای نوار ذخیره و به‌روزرسانی می‌کند. حالت فعلی ماشین دوبعدی نیز در کنترل متناهی ماشین تک‌نواره ذخیره می‌شود.
    
    \item {شبیه‌سازی هر گام محاسباتی:} 

    برای شبیه‌سازی یک حرکت، ماشین یک‌بعدی ابتدا نوار خود را از ابتدا تا انتها جستجو می‌کند تا مختصات فعلی $(x_c, y_c)$ را پیدا کند. 
    اگر چنین رکوردی پیدا شد، نماد آن را می‌خواند. اگر پیدا نشد، نماد آن خانه را خالی فرض می‌کند.
    سپس با توجه به تابع انتقال ماشین دوبعدی، نماد جدید را می‌نویسد (اگر رکوردی وجود داشت آن را بازنویسی می‌کند و اگر نداشت، یک رکورد جدید به انتهای نوار اضافه می‌کند)، حالت جدید را در کنترل متناهی خود ثبت می‌کند و مختصات فعلی $(x_c, y_c)$ را بر اساس جهت حرکت به‌روزرسانی می‌کند (مثلا برای حرکت به بالا، $y_c$ را یک واحد افزایش و برای حرکت به چپ، $x_c$ را یک واحد کاهش می‌دهد).
    
    \item {مدیریت ورودی و خروجی:}

     در ابتدا، ماشین یک‌بعدی رشته ورودی را می‌خواند و آن را به فرمت $(0,0,w_1)\#(1,0,w_2)\dots$ روی نوار خود قالب‌بندی می‌کند. 
    هنگامی که ماشین دوبعدی متوقف می‌شود، طبق صورت سوال باید فقط خروجی موجود در سطر فعلی هد در نظر گرفته شود. بنابراین ماشین یک‌بعدی تمام رکوردهایی که دارای مختصات $y$ برابر با $y_c$ (موقعیت نهایی هد) هستند را پیدا کرده، نمادهای آن‌ها را به ترتیب $x$ در ابتدای نوار مرتب می‌کند و بقیه محتوای نوار را پاک می‌کند تا خروجی نهایی دقیقا منطبق با رفتار ماشین دوبعدی تولید شود.
\end{itemize}

از آنجا که هر ماشین تورینگ استاندارد می‌تواند توسط ماشین دوبعدی شبیه‌سازی شود و ماشین دوبعدی نیز می‌تواند توسط ماشین استاندارد شبیه‌سازی شود، هر دو مدل از نظر قدرت محاسباتی معادل یکدیگر هستند.
